\paragraph{Los vectores y sus tres operaciones.} Las matrices del proceso se aplican siempre sobre \emph{vectores}, y es con los vectores con los que se hacen las tres operaciones del curso: la norma $\lVert\vect{a}\rVert$, el producto punto $\vect{a}\cdot\vect{b}$ y el producto cruz $\vect{a}\times\vect{b}$. Los párrafos siguientes siguen, con los números del casco, cada vector que aparece en este paso y dicen qué significa cada operación.

\paragraph{Del píxel al rayo.} Un píxel es el vector $\vect{p}=(u,v,1)^{T}$ (coordenadas homogéneas). Al multiplicarlo por $K^{-1}$ se obtiene
el vector $\vect{r}=K^{-1}\vect{p}=((u-c_x)/f,\ (v-c_y)/f,\ 1)^{T}$, que apunta desde la cámara hacia el punto del mundo que vio ese píxel. Para
el píxel A, $(u,v)=(488,576)$, por ejemplo,
\begin{equation}
\vect{r}_A=\begin{pmatrix}\dfrac{488-480}{753{,}449}\\[2mm]\dfrac{576-640}{753{,}449}\\[1mm]1\end{pmatrix}=\begin{pmatrix} 0{,}0106 \\ -0{,}0849 \\ 1 \end{pmatrix} .
\end{equation}
La \emph{norma} de $\vect{r}$ es la longitud de ese vector, $\lVert\vect{r}_A\rVert=\sqrt{(0{,}0106)^2+(-0{,}0849)^2+1^2}=1{,}0037$, y el ángulo
que forma con el eje óptico $\vect{e}_z=(0,0,1)^{T}$ sale del producto punto:
\begin{equation}
\cos\theta_A=\frac{\vect{r}_A\cdot\vect{e}_z}{\lVert\vect{r}_A\rVert\,\lVert\vect{e}_z\rVert}=\frac{1}{1{,}0037}=0{,}9964,\qquad \theta_A=4{,}89^\circ .
\end{equation}
Cuanto más lejos del centro está el píxel, más inclinado es su rayo (cuadro~\ref{tab:rayos-casco} y figura~\ref{fig:vec-rayos}).

\begin{table}[htbp]
\caption{Los tres píxeles del casco como vectores: píxel $\vect{p}$, rayo $\vect{r}=K^{-1}\vect{p}$, su norma, su ángulo con el eje óptico y la profundidad $Z$ (mm)}
\label{tab:rayos-casco}
\centering\small
\begin{tabular}{@{}clcccc@{}}
\toprule
Píxel & $\vect{p}=(u,v,1)$ & $\vect{r}=K^{-1}\vect{p}$ & $\lVert\vect{r}\rVert$ & $\theta$ & $Z$ \\
\midrule
A & $(488,\,576,\,1)$ & $(0{,}0106,\,-0{,}0849,\,1)$ & $1{,}0037$ & $4{,}89^\circ$ & $300{,}9$ \\
B & $(450,\,300,\,1)$ & $(-0{,}0398,\,-0{,}4513,\,1)$ & $1{,}0978$ & $24{,}37^\circ$ & $307{,}9$ \\
C & $(221,\,576,\,1)$ & $(-0{,}3438,\,-0{,}0849,\,1)$ & $1{,}0608$ & $19{,}50^\circ$ & $379{,}1$ \\
\bottomrule
\end{tabular}
\end{table}

\paragraph{Del rayo al punto.} Multiplicar el rayo por la profundidad lo estira hasta el punto del casco: $\vect{P}=Z\,\vect{r}$. Es una
multiplicación de un vector por un escalar, y el resultado conserva la dirección. Para A, con $Z=300{,}87$ mm,
\begin{equation}
\vect{P}_A=300{,}87\cdot\begin{pmatrix} 0{,}0106 \\ -0{,}0849 \\ 1 \end{pmatrix}=\begin{pmatrix} 3{,}19 \\ -25{,}56 \\ 300{,}87 \end{pmatrix}\ \text{mm},
\end{equation}
y su norma, $\lVert\vect{P}_A\rVert=301{,}97$ mm, es la distancia real de la cámara a ese punto (mayor que $Z$ porque el rayo está inclinado:
$\lVert\vect{P}\rVert=Z\lVert\vect{r}\rVert$).

\paragraph{Producto punto entre dos rayos.} El ángulo entre los rayos de dos píxeles se calcula igual. Para A y C,
\begin{equation}
\vect{r}_A\cdot\vect{r}_C=(0{,}0106)(-0{,}3438)+(-0{,}0849)(-0{,}0849)+1\cdot1=1{,}0036,
\end{equation}
\begin{equation}
\cos\theta_{AC}=\frac{1{,}0036}{(1{,}0037)(1{,}0608)}=0{,}9426,
\end{equation}
es decir $\theta_{AC}=19{,}51^\circ$. En la figura~\ref{fig:vec-rayos} están los ángulos de los tres píxeles con el eje óptico y entre sí.

\begin{figure}[htbp]
\centering
\includegraphics[width=0.97\linewidth]{vec_rayos.jpg}
\caption{Izquierda: los rayos $\vect{r}_A$, $\vect{r}_B$ y $\vect{r}_C$ (con componente $z=1$) cruzan el plano de la imagen, y sus puntos
$\vect{P}=Z\vect{r}$ están más lejos sobre el mismo rayo. Derecha: ángulos de cada rayo con el eje óptico y entre rayos, calculados con el
producto punto.}
\label{fig:vec-rayos}
\end{figure}
